\chapter{Dirichlet Characters and \texorpdfstring{$L$}{L}-Functions}\label{ch:characters}

\begin{center}\itshape
Applying Fourier series to number theory, Dirichlet has recently found results touching the utmost of human acumen.\\
\upshape --- Carl Gustav Jacob Jacobi, letter to his brother Moritz Hermann Jacobi, 10 September 1838
\end{center}

Start with a series you can sum. Leibniz's series is
\[
1-\frac13+\frac15-\frac17+\cdots=\frac\pi4.
\]
The number $\frac\pi4$ is about $0.785$, and it is not $0$. The signs in that series are a function of $n$. Write $\chi_{-4}(n)=0,1,0,-1$ according as $n\equiv0,1,2,3\pmod4$. Then the series is $L(1,\chi_{-4})=\sum_{n\ge1}\chi_{-4}(n)n^{-1}$.

Dirichlet's question was about primes. The progression $5,15,25,35,\dots$ has difference $10$ and first term $5$. Those two integers share the factor $5$, and every later term is divisible by $5$. Move the first term to $2$. The progression $2,7,12,17,22,27,32,37,\dots$ has difference $5$, and $\gcd(2,5)=1$. The primes $2,7,17,37$ sit in it. Does the list go on forever?

Euclid's argument answers that question for a few progressions, then stops. On 27 July 1837 Dirichlet, then 32 and lecturing in Berlin, read the Academy a proof that works for every progression whose first term and difference are coprime. The proof splits the primes by characters. Each character has a series $L(s,\chi)$. The principal character produces a pole at $s=1$. The proof survives only if no other character has $L(1,\chi)=0$.

A primitive character has a Gauss sum of modulus $\sqrt q$. Chapter~\ref{ch:sums} uses that sum to bound character sums. The finite formula for $L(1,\chi)$ carries the class numbers of quadratic fields. The fact proved here is $L(1,\chi)\ne0$ when $\chi\ne\chi_0$. Chapter~\ref{ch:dirichlet} uses it to share the primes among the classes.

\begin{think}
For $\chi_{-4}$ as above, $L(1,\chi_{-4})=\frac\pi4$. Formally,
$\log L(s,\chi_{-4})\approx\sum_{p\equiv1\,(4)}p^{-s}-\sum_{p\equiv3\,(4)}p^{-s}$ as $s\to1^+$. What does $L(1,\chi_{-4})\ne0$ say about the primes in the two classes?
What would $L(1,\chi_{-4})=0$ have meant?
\end{think}

\section{Characters}

Fix the modulus $4$ before the general definition. The units modulo $4$ are $1$ and $3$. There are two homomorphisms from that group to the unit circle. One sends both units to $1$. The other sends $3$ to $-1$. Extend each by the value $0$ at even $n$. You get two functions on $\Z$: the principal character $\chi_0$, equal to $1$ on odd $n$ and $0$ on even $n$, and $\chi_{-4}$ as in the opening. Now check one period of $\chi_{-4}$:
\[
\chi_{-4}(1)+\chi_{-4}(2)+\chi_{-4}(3)+\chi_{-4}(4)=1+0-1+0=0.
\]
The partial sums $1,1,0,0,1,1,0,\dots$ never leave the set $\{0,1\}$.

A \emph{Dirichlet character modulo $q$} (Volume~3, Chapter~11) is the same construction for a general $q\ge1$. It is a function $\chi:\Z\to\C$ with $\chi(n)=\psi(n\bmod q)$ when $\gcd(n,q)=1$, for a group homomorphism $\psi:(\Z/q\Z)^*\to\C^*$, and with $\chi(n)=0$ when $\gcd(n,q)>1$. So $\chi$ is $q$-periodic and completely multiplicative. Its non-zero values are $\varphi(q)$-th roots of unity. There are $\varphi(q)$ characters modulo $q$. They form a group $\widehat G_q$ under multiplication. The \emph{principal} character $\chi_0$ is the indicator of the integers prime to $q$. The inverse of $\chi$ is $\overline\chi$. A character is \emph{real} when its values are real, which is the same condition as $\chi^2=\chi_0$. Otherwise it is \emph{complex}. It is \emph{even} or \emph{odd} according as $\chi(-1)=1$ or $\chi(-1)=-1$.

The sum over one period is the property the later estimates use.

\begin{proposition}[orthogonality]\label{prop:orth}
\begin{enumerate}[label=(\alph*)]
\item For $\chi\ne\chi_0$, $\sum_{n=1}^q\chi(n)=0$, and $\bigl|\sum_{n\le N}\chi(n)\bigr|\le\varphi(q)$ for every $N\ge1$.
\item For $a$ prime to $q$ and every $n$, $\frac1{\varphi(q)}\sum_\chi\overline{\chi(a)}\chi(n)$ is $1$ if $n\equiv a\pmod q$, and $0$ otherwise.
\end{enumerate}
\end{proposition}
\begin{proof}
(a) Since $\chi$ is not principal, some $b$ prime to $q$ has $\chi(b)\ne1$. As $n$ runs through a complete set of residues modulo $q$, so does $bn$. Therefore
\[
\chi(b)\sum_{n=1}^q\chi(n)=\sum_{n=1}^q\chi(bn)=\sum_{n=1}^q\chi(n).
\]
The factor $\chi(b)-1$ is not zero, so the sum over a period is $0$.

A long sum now collapses. The sum of $\chi$ over each complete block of $q$ consecutive integers is $0$, by periodicity. The sum up to $N$ is the sum over the incomplete last block. That block has length less than $q$. At most $\varphi(q)$ of its terms are prime to $q$, and each of those has absolute value $1$. The bound follows.

(b) is Volume~3, Proposition~11.15.
\end{proof}

Part (b) is the chord split into notes. The indicator of the class $a$ is the average of the characters against $\overline{\chi(a)}$. Each character is multiplicative. The indicator itself is not.

\begin{example}\label{ex:smallchars}
Modulo $4$ there are two characters, $\chi_0$ and $\chi_4=\chi_{-4}$ above. The second is odd, because $\chi_{-4}(-1)=\chi_{-4}(3)=-1$. Modulo $3$ there are again two: $\chi_0$, and $\chi_{-3}(n)=0,1,-1$ for $n\equiv0,1,2$. The character $\chi_{-3}$ is odd. Modulo $5$ the group $(\Z/5\Z)^*$ is cyclic of order $4$, generated by $2$. A character is fixed by the value $\chi(2)\in\{1,i,-1,-i\}$. The choice $\chi(2)=-1$ gives the Legendre symbol $\chi_5(n)=\bigl(\frac n5\bigr)$, which is even. The choices $\chi(2)=\pm i$ give two complex odd characters. Call the character with value $i$ at $2$ by the name $\psi$, so $\overline\psi$ is the character with value $-i$. Figure~\ref{fig:mod5grid} is the whole table, including the column $n\equiv0$, where every character vanishes. For an odd prime $p$, the Legendre symbol $\bigl(\frac\cdot p\bigr)$ is the unique character of order $2$ modulo $p$ (Volume~3, Chapter~11). It is odd exactly when $p\equiv3\pmod4$, because $\bigl(\frac{-1}p\bigr)=(-1)^{(p-1)/2}$.
\end{example}

\begin{figure}[ht]
\centering
\fig{ch03-new-mod5}
\caption{The four characters modulo $5$. Columns are the residues $n$. Coral bars mark the complex rows: $\psi$ with $\psi(2)=i$, and its conjugate $\overline\psi$. The Legendre symbol is the real row with value $-1$ at $n=2$. Every character is $0$ at multiples of $5$.}
\label{fig:mod5grid}
\end{figure}

Proposition~\ref{prop:orth}(a) bounds the partial sums of a non-principal character by $\varphi(q)$. The sums can still be large beside $\sqrt q$, and the signs are not independent. Figure~\ref{fig:charsums} shows a steady bias. The remark on class numbers, at the end of the chapter, names the height of that bias.

\begin{figure}[ht]
\centering
\fig{ch03-charsums}
\caption{The partial sums $S(N)=\sum_{n\le N}\bigl(\frac n{1019}\bigr)$. Since $1019\equiv3\pmod4$, the symbol is odd and $S(1019-1-N)=S(N)$. The walk stays
almost entirely above $0$, and reaches $39=3\cdot13$ at $N=509$, where $13$ is the class number of $\Q(\sqrt{-1019})$
(Exercise~\ref{ex:residuesum} and the remark on class numbers).}
\label{fig:charsums}
\end{figure}

\section{Primitive characters}

Some characters are written at a larger modulus than they need. Modulo $12$, keep $\chi_{-4}$ on the integers prime to $12$ and set the value $0$ on the multiples of $3$. On $1,5,7,11$ the values are $1,1,-1,-1$, because those residues are $1,1,3,3$ modulo $4$. This function is a character modulo $12$. It is $\chi_{-4}$ with the multiples of $3$ removed. Its natural modulus is $4$.

To measure that, for $d\mid q$ let $U_d=\{n\in(\Z/q\Z)^*:n\equiv1\pmod d\}$. This is a subgroup of the units.

\begin{lemma}\label{lem:lift}
Let $d\mid q$ and $m$ prime to $d$. There is an $n\equiv m\pmod d$ with $\gcd(n,q)=1$.
\end{lemma}
\begin{proof}
Let $t$ be the product of the primes that divide $q$ and do not divide $m$, and set $n=m+td$. Take a prime $p$ that divides $q$. If $p$ divides $m$, then $p$ does not divide $d$, because $\gcd(m,d)=1$, and $p$ does not divide $t$ by the way $t$ was built. So $p$ does not divide $td$, and $p$ does not divide $n$. If $p$ does not divide $m$, then $p$ divides $t$, so again $p$ does not divide $n$. Every prime divisor of $q$ misses $n$.
\end{proof}

The conductor is the smallest modulus that still carries the character.

\begin{proposition}\label{prop:conductor}
Let $\chi$ be a character modulo $q$, and let $D$ be the set of $d\mid q$ such that $\chi=1$ on $U_d$.
\begin{enumerate}[label=(\alph*)]
\item If $d_1,d_2\in D$ then $\gcd(d_1,d_2)\in D$. So the least element $d^*$ of $D$, the \emph{conductor} of $\chi$, divides every element of $D$.
\item For $d\in D$ there is a unique character $\psi$ modulo $d$ with $\chi(n)=\psi(n)$ for all $n$ prime to $q$; we say $\psi$ \emph{induces} $\chi$. Conversely,
if a character modulo $d\mid q$ induces $\chi$, then $d\in D$.
\item Call $\chi$ \emph{primitive} if $d^*=q$. Exactly one primitive character induces $\chi$: the character $\chi^*$ modulo $d^*$ given by (b).
\end{enumerate}
\end{proposition}
\begin{proof}
(a) Write $q=\prod p^{e_p}$. By the Chinese remainder theorem, $(\Z/q\Z)^*\cong\prod_p(\Z/p^{e_p}\Z)^*$. The condition $n\equiv1\pmod d$ means $n\equiv1\pmod{p^{v_p(d)}}$ for every $p$. So $U_d=\prod_p V_p(v_p(d))$, where $V_p(f)$ is the subgroup of residues $\equiv1\pmod{p^f}$ inside $(\Z/p^{e_p}\Z)^*$. These subgroups shrink as $f$ grows. Hence
\[
U_{d_1}U_{d_2}=\prod_p V_p\bigl(\min(v_p(d_1),v_p(d_2))\bigr)=U_{\gcd(d_1,d_2)}.
\]
If $\chi$ is trivial on $U_{d_1}$ and on $U_{d_2}$, it is trivial on the product. If $d\in D$, then $\gcd(d,d^*)\in D$. That gcd is at most $d^*$, so it equals $d^*$, which means $d^*\mid d$.

(b) Take $m$ prime to $d$. Lemma~\ref{lem:lift} supplies an $n\equiv m\pmod d$ with $\gcd(n,q)=1$. Set $\psi(m)=\chi(n)$. If $n'$ is another such lift, then $n'\equiv n\pmod d$, so $n'n^{-1}\in U_d$, the inverse being taken modulo $q$. Thus $\chi(n')=\chi(n)$. The value $\psi(m)$ does not depend on the lift. It is multiplicative, because a product of lifts is a lift of the product. Extend $\psi$ by $0$ off the residues prime to $d$. Then $\psi$ is a character modulo $d$, and it agrees with $\chi$ on integers prime to $q$. It is the only such character: every residue prime to $d$ contains an integer prime to $q$, again by the lemma. Conversely, if a character $\psi$ modulo $d$ induces $\chi$, then for $n\in U_d$ one has $\chi(n)=\psi(n)=\psi(1)=1$, so $d\in D$.

(c) Let $\chi^*$ be the character modulo $d^*$ that induces $\chi$. Suppose $\chi^*$ were trivial on its own subgroup $U'_{d'}$ for some $d'\mid d^*$ with $d'<d^*$. For $n\in U_{d'}$ one would have $\chi(n)=\chi^*(n)=1$, so $d'\in D$, against the minimality of $d^*$. Thus $\chi^*$ is primitive. Now let $\psi$ modulo $d$ be primitive and induce $\chi$. Part (b) puts $d$ in $D$, so $d^*\mid d$ by (a). Take $m$ prime to $d$ with $m\equiv1\pmod{d^*}$. Lift $m$ to an $n$ prime to $q$, by Lemma~\ref{lem:lift}. Then $n\in U_{d^*}$, so $\psi(m)=\chi(n)=1$. Thus $\psi$ is trivial on the subgroup attached to $d^*$ modulo $d$. Primitivity forces $d=d^*$. Uniqueness in (b) forces $\psi=\chi^*$.
\end{proof}

For example, $\chi_{-4}$, $\chi_{-3}$ and $\chi_5$ are primitive, and the conductor of $\chi_0$ is $1$. For a prime modulus every non-principal character is primitive. The only proper divisor of the modulus is $1$, and $U_1$ is the whole unit group, on which a non-principal character is not trivial.

\section{Gauss sums}

Characters see multiplication. The exponentials $e(a/q)$ see addition. A Gauss sum is the meeting place. It returns in the formula for $L(1,\chi)$ below, in the sieve of Chapter~\ref{ch:sieves}, and in the P\'olya--Vinogradov inequality of Chapter~\ref{ch:sums}. Write $e(x)=e^{2\pi ix}$. The \emph{Gauss sum} of a character $\chi$ modulo $q$ is
\[
\tau(\chi)=\sum_{a=1}^q\chi(a)e\Bigl(\frac aq\Bigr).
\]
For the Legendre symbol modulo an odd prime $p$, this is the quadratic Gauss sum of Volume~5, Proposition~7.13. Gauss found its value, including the sign (Volume~5, Theorem~8.14). This chapter uses the modulus $\sqrt q$ for a general primitive character, and leaves that sign where Volume~5 proved it.

\begin{theorem}\label{thm:gauss}
Let $\chi$ be a primitive character modulo $q$. Then for every integer $n$,
\[
\sum_{a=1}^q\chi(a)e\Bigl(\frac{an}q\Bigr)=\overline{\chi(n)}\,\tau(\chi),
\]
and $|\tau(\chi)|=\sqrt q$.
\end{theorem}
\begin{proof}
First take $\gcd(n,q)=1$. Substitute $a=n^{-1}b$, with the inverse modulo $q$. As $a$ runs through the residues, so does $b$. Since $\chi(n)\chi(n^{-1})=1$ and $|\chi(n)|=1$, one has $\chi(a)=\overline{\chi(n)}\,\chi(b)$. The sum becomes $\overline{\chi(n)}\sum_b\chi(b)e(b/q)=\overline{\chi(n)}\,\tau(\chi)$.

Now take $g=\gcd(n,q)>1$. The right side is $0$, because $\chi(n)=0$. Write $T$ for the left side, $q=gq_1$ and $n=gn_1$. Then $e(an/q)=e(an_1/q_1)$ depends only on $a$ modulo $q_1$. Here $q_1$ is a proper divisor of $q$. Primitivity says $q_1\notin D$, so some $c\in U_{q_1}$ has $\chi(c)\ne1$. As $a$ runs through the residues modulo $q$, so does $ca$, and $ca\equiv a\pmod{q_1}$. Therefore
\[
T=\sum_a\chi(ca)\,e\Bigl(\frac{an_1}{q_1}\Bigr)=\chi(c)\,T.
\]
Since $\chi(c)\ne1$, one gets $T=0$.

For the modulus, sum the squared absolute values of both sides over $n=1,\dots,q$, and call the common total $T$. On the right, $|\overline{\chi(n)}|$ is $1$ for the $\varphi(q)$ residues prime to $q$ and is $0$ otherwise, so $T=\varphi(q)\,|\tau(\chi)|^2$. On the left, expand:
\[
\begin{aligned}T&=\sum_{a,b}\chi(a)\overline{\chi(b)}\sum_{n=1}^qe\Bigl(\frac{(a-b)n}q\Bigr)\eqbr=q\sum_a|\chi(a)|^2=q\varphi(q),\end{aligned}
\]
because the inner geometric sum equals $q$ when $a\equiv b\pmod q$ and equals $0$ otherwise. Hence $|\tau(\chi)|^2=q$.
\end{proof}

So a primitive character is, up to the constant $\tau(\chi)$, its own Fourier transform on $\Z/q\Z$. Every frequency has size $\sqrt q$. This is the character form of Volume~5, Proposition~7.13, and it is the input to the P\'olya--Vinogradov inequality in Chapter~\ref{ch:sums}. Here it produces a finite formula for $L(1,\chi)$.

\section{\texorpdfstring{$L$}{L}-functions}

The \emph{Dirichlet $L$-function} of $\chi$ is $L(s,\chi)=\sum_{n\ge1}\chi(n)n^{-s}$. We write $s=\sigma+it$.

\begin{theorem}\label{thm:lbasic}
Let $\chi$ be a character modulo $q$.
\begin{enumerate}[label=(\alph*)]
\item For $\sigma>1$ the series converges absolutely, $L(s,\chi)=\prod_p\bigl(1-\chi(p)p^{-s}\bigr)^{-1}$, and
\[
L(s,\chi)=\exp\Bigl(\sum_p\sum_{k\ge1}\frac{\chi(p)^k}{kp^{ks}}\Bigr).
\]
In particular $L(s,\chi)\ne0$ for $\sigma>1$.
\item If $\chi\ne\chi_0$, the series converges on $\sigma>0$ and $L(s,\chi)$ is holomorphic there.
\item $L(s,\chi_0)=\zeta(s)\prod_{p\mid q}(1-p^{-s})$ for $\sigma>1$, so $L(s,\chi_0)-\frac{\varphi(q)/q}{s-1}$ extends holomorphically to $\sigma>0$.
\end{enumerate}
\end{theorem}
\begin{proof}
(a) The comparison $|\chi(n)n^{-s}|\le n^{-\sigma}$ gives absolute convergence for $\sigma>1$. The function $\chi$ is completely multiplicative, so Volume~5, Theorem~10.3 supplies the Euler product. The local factor is the geometric series $\sum_{k\ge0}\chi(p)^kp^{-ks}=(1-\chi(p)p^{-s})^{-1}$.

For the exponential, take $|w|<1$. The function $E(w)=\exp\bigl(\sum_{k\ge1}w^k/k\bigr)$ is holomorphic on the disc, $E(0)=1$, and $E'=E/(1-w)$. Thus $((1-w)E)'=0$, so $(1-w)E$ is constant, and $E(w)=1/(1-w)$. Set $w=\chi(p)p^{-s}$. The partial Euler product over $p\le y$ equals $\exp\bigl(\sum_{p\le y}\sum_{k\ge1}\chi(p)^k/(kp^{ks})\bigr)$. The double series converges absolutely, because its terms are bounded by those of $\sum_p\sum_{k\ge1}p^{-k\sigma}$, which converges for $\sigma>1$. Let $y\to\infty$ and use the continuity of the exponential. An exponential never vanishes, so $L(s,\chi)\ne0$ for $\sigma>1$.

(b) Proposition~\ref{prop:orth}(a) bounds the partial sums of $\chi$. Volume~5, Proposition~10.1(b) turns a bounded partial-sum sequence into a holomorphic Dirichlet series on $\sigma>0$.

(c) The value $\chi_0(n)$ is $1$ exactly when no prime divisor of $q$ divides $n$. The Euler product of $L(s,\chi_0)$ is the Euler product of $\zeta$, with the factors at $p\mid q$ removed. By Volume~5, Proposition~10.6, $\zeta(s)-1/(s-1)$ is holomorphic on $\sigma>0$. The product $P(s)=\prod_{p\mid q}(1-p^{-s})$ is entire. Therefore
\[
L(s,\chi_0)-\frac{P(1)}{s-1}=P(s)\Bigl(\zeta(s)-\frac1{s-1}\Bigr)+\frac{P(s)-P(1)}{s-1}
\]
extends holomorphically to $\sigma>0$. And $P(1)=\prod_{p\mid q}(1-1/p)=\varphi(q)/q$.
\end{proof}

If $\chi^*$ modulo $d$ induces $\chi$ modulo $q$, the Euler products differ only at the primes that divide $q$:
\[
L(s,\chi)=L(s,\chi^*)\prod_{p\mid q}\bigl(1-\chi^*(p)p^{-s}\bigr).
\]
Questions about $L(s,\chi)$ near $s=1$ reduce to the primitive character $\chi^*$. The extra factors are finite in number and do not vanish at $s=1$.

Figure~\ref{fig:lsigmachi4} is a numerical check for the one series from the opening, before the proof that the value at $\sigma=1$ cannot be zero. The five heights were computed from the paired series $\sum_{k\ge0}\bigl((4k+1)^{-\sigma}-(4k+3)^{-\sigma}\bigr)$. At $\sigma=1$ the sum is $\pi/4$.

\begin{figure}[ht]
\centering
\fig{ch03-new-lvals}
\caption{Five computed values of $L(\sigma,\chi_{-4})$, rounded to three decimals from $0.6677$, $0.7854$, $0.8645$, $0.9160$ and $0.9689$. The coral bar is $\sigma=1$, where the value is $\pi/4$. The vertical scale starts at $0.60$. Every bar is positive.}
\label{fig:lsigmachi4}
\end{figure}

\section{The non-vanishing of \texorpdfstring{$L(1,\chi)$}{L(1,chi)}}\label{sec:siegel}

The next theorem is the step Dirichlet's argument cannot spare. One zero of a non-principal $L$-function at $s=1$ would cancel the pole that produces the primes. For a complex character the product of all the $L$-functions rules that out. For a real character the same product has only one suspicious factor, and the proof uses a series with non-negative coefficients. Figure~\ref{fig:lreal} draws three real examples. All three are positive at $\sigma=1$.

\begin{figure}[ht]
\centering
\fig{ch03-lreal}
\caption{$L(\sigma,\chi)$ for $0<\sigma\le3$ and the real characters $\chi_{-3}$, $\chi_{-4}$, $\chi_5$. At $\sigma=1$ the values are $\frac\pi{3\sqrt3}=0.6046$, $\frac\pi4=0.7854$ and
$\frac2{\sqrt5}\log\frac{1+\sqrt5}2=0.4304$ (Theorem~\ref{thm:l1formula}). As $\sigma\to\infty$ all tend to $\chi(1)=1$. The even character $\chi_5$ tends to $0$ as $\sigma\to0$.}
\label{fig:lreal}
\end{figure}

\begin{lemma}\label{lem:prodchars}
For real $\sigma>1$, $\prod_{\chi\bmod q}L(\sigma,\chi)\ge1$.
\end{lemma}
\begin{proof}
Theorem~\ref{thm:lbasic}(a) writes the product as $\exp(S)$, where
\[
\begin{aligned}S&=\sum_p\sum_{k\ge1}\frac1{kp^{k\sigma}}\sum_\chi\chi(p^k)\eqbr=\varphi(q)\sum_{p^k\equiv1\,(q)}\frac1{kp^{k\sigma}}\ge0,\end{aligned}
\]
by Proposition~\ref{prop:orth}(b) with $a=1$. The sums may be exchanged: the double series converges absolutely. So $S$ is real and non-negative, and $\exp(S)\ge1$.
\end{proof}

The inner sum over characters keeps only the prime powers that are $1$ modulo $q$. For $q=4$ those are the powers $p^k\equiv1\pmod4$. Their coefficients are positive, which is why the product cannot tend to $0$.

\begin{theorem}[Dirichlet, 1837]\label{thm:lnonzero}
$L(1,\chi)\ne0$ for every character $\chi\ne\chi_0$.
\end{theorem}
\begin{proof}
\emph{Complex characters.} Suppose $L(1,\chi)=0$ with $\chi\ne\overline\chi$. Conjugating the series, which converges at $s=1$, gives $L(1,\overline\chi)=\overline{L(1,\chi)}=0$. The conjugate is a second factor of the product in Lemma~\ref{lem:prodchars}. Let $\sigma\to1^+$. Comparing $\sum n^{-\sigma}$ with $1+\int_1^\infty t^{-\sigma}\,\dd t$ gives $0<L(\sigma,\chi_0)\le\zeta(\sigma)\le\sigma/(\sigma-1)$. Each of $L(\sigma,\chi)$ and $L(\sigma,\overline\chi)$ is holomorphic near $1$ and has a zero there, so each is $O(\sigma-1)$. Every remaining factor is continuous at $1$, hence bounded. The product is $O\bigl((\sigma-1)^{-1}(\sigma-1)^2\bigr)$ and tends to $0$. That contradicts Lemma~\ref{lem:prodchars}.

\emph{Real characters.} Suppose $L(1,\chi)=0$ with $\chi$ real and $\chi\ne\chi_0$. Let $r=\mathbf1*\chi$, so $r(n)=\sum_{d\mid n}\chi(d)$. The function $r$ is multiplicative. At a prime power, $r(p^k)=\sum_{j=0}^k\chi(p)^j$. This equals $k+1$ if $\chi(p)=1$, equals $1$ if $\chi(p)=0$, and equals $1$ or $0$ according as $k$ is even or odd if $\chi(p)=-1$. Hence $r(n)\ge0$ for every $n$, and $r(m^2)\ge1$. By Volume~5, Proposition~10.2(a), $F(s)=\zeta(s)L(s,\chi)=\sum r(n)n^{-s}$ for $\sigma>1$.

The difference $\zeta(s)-1/(s-1)$ is holomorphic on $\sigma>0$. So is $L(s,\chi)/(s-1)$, extended by the value $L'(1,\chi)$ at $s=1$, because $L(1,\chi)=0$. Therefore
\[
F(s)=\Bigl(\zeta(s)-\frac1{s-1}\Bigr)L(s,\chi)+\frac{L(s,\chi)}{s-1}
\]
extends holomorphically to $\sigma>0$. That half-plane contains the disc $|s-2|<2$. By Volume~5, Theorem~9.11, $F$ equals its Taylor series at $2$ throughout the disc. Differentiating the Dirichlet series $k$ times, by Volume~5, Proposition~10.1(a), gives $F^{(k)}(2)=\sum_n r(n)(-\log n)^k n^{-2}$. At the point $s=\frac12$, which lies in the disc,
\[
\begin{aligned}F\Bigl(\frac12\Bigr)&=\sum_{k\ge0}\frac{F^{(k)}(2)}{k!}\Bigl(-\frac32\Bigr)^k\eqbr=\sum_{k\ge0}\sum_{n\ge1}\frac{r(n)}{n^2}\frac{(\frac32\log n)^k}{k!}\eqbr=\sum_{n\ge1}\frac{r(n)}{n^2}e^{\frac32\log n}\eqbr=\sum_{n\ge1}\frac{r(n)}{\sqrt n}.\end{aligned}
\]
Every term $r(n)n^{-2}(\frac32\log n)^k/k!$ is non-negative, so the two sums may be exchanged. The series $\sum r(n)n^{-1/2}$ therefore converges. It is at least $\sum_m r(m^2)m^{-1}\ge\sum_m m^{-1}$, and the harmonic series diverges. This contradiction proves the theorem.
\end{proof}

\begin{beauty}
The two halves use positivity in different places. For complex $\chi$, the product of all the $L$-functions is the generating function of the prime powers $p^k\equiv1\pmod q$. Those coefficients are non-negative. Two zeros would beat the single pole of $L(s,\chi_0)$, and the product would tend to $0$. For real $\chi$ only one zero is on offer. The positive object is then $\zeta(s)L(s,\chi)$, the zeta function of the quadratic field attached to $\chi$. Its coefficient $r(n)$ counts ideals of norm $n$. A Dirichlet series with non-negative coefficients cannot be holomorphic past the point where it diverges. The argument above is that fact for this one series, which is the special case of Landau's theorem used here. The squares force a divergence at $\frac12$. The real case is the open one: how small $L(1,\chi)$ can be, for real $\chi$, is still unsettled. A real zero very close to $s=1$ is called a Siegel zero.
\end{beauty}

\section{The value of \texorpdfstring{$L(1,\chi)$}{L(1,chi)}}

Non-vanishing says the number is not zero. Dirichlet computed it. For $\chi_{-4}$ the value is $\pi/4$, which is Leibniz's series again. The input is the sum of $z^n/n$ on the unit circle.

On 10 September 1838 Jacobi wrote to his brother Moritz Hermann Jacobi. The sentence at the head of the chapter is Elstrodt's English for that letter. The series in the next lemma is one of the Fourier series in that work. The class-number sums Jacobi was praising use the same value $L(1,\chi)$.

\begin{lemma}\label{lem:logseries}
For $0<x<2\pi$, $\displaystyle\sum_{n\ge1}\frac{e^{inx}}n=-\log\Bigl(2\sin\frac x2\Bigr)+i\,\frac{\pi-x}2$.
\end{lemma}
\begin{proof}
Let $z=e^{ix}\ne1$. For $t\in[0,1]$, the point $1-zt$ is not $0$: the equality $|zt|=1$ would force $t=1$, and $z\ne1$. By compactness, $|1-zt|\ge c>0$ on $[0,1]$. Then
\[
\sum_{n=1}^N\frac{z^n}n=\int_0^1\sum_{n=1}^Nz^nt^{n-1}\dd t=\int_0^1\frac{z-z^{N+1}t^N}{1-zt}\dd t,
\]
and the piece carrying $t^N$ has absolute value at most $c^{-1}\int_0^1 t^N\,\dd t$, which tends to $0$. The series converges to $\int_0^1 z/(1-zt)\,\dd t$.

On $[0,1[$ one has $\operatorname{Re}(1-zt)=1-t\cos x>0$. The principal logarithm $-\operatorname{Log}(1-zt)$ is a $C^1$ primitive of $z/(1-zt)$ on $[0,1]$, so the integral equals $-\operatorname{Log}(1-z)$. Finally
\[
1-e^{ix}=-2i\sin\frac x2\,e^{ix/2}=2\sin\frac x2\,e^{i(x-\pi)/2}.
\]
Here $2\sin(x/2)>0$, and $(x-\pi)/2$ lies in $\bigl]-\pi/2,\pi/2\bigr[$, so this is the principal argument. Therefore $\operatorname{Log}(1-z)=\log\bigl(2\sin(x/2)\bigr)+i(x-\pi)/2$.
\end{proof}

The imaginary part is $\sum(\sin nx)/n=(\pi-x)/2$, which is the Fourier series from Volume~5, Chapter~6. The real part is $\sum(\cos nx)/n=-\log\bigl(2\sin(x/2)\bigr)$.

\begin{theorem}[Dirichlet's formula]\label{thm:l1formula}
Let $\chi$ be a primitive character modulo $q>1$. Then
\[
\begin{gathered}L(1,\chi)=-\frac{i\pi}{q\,\tau(\overline\chi)}\sum_{a=1}^{q-1}\overline\chi(a)\,a\quad(\chi\text{ odd}),\\
L(1,\chi)=-\frac1{\tau(\overline\chi)}\sum_{a=1}^{q-1}\overline\chi(a)\log\sin\frac{\pi a}q\quad(\chi\text{ even}).\end{gathered}
\]
\end{theorem}
\begin{proof}
The conjugate $\overline\chi$ is primitive, and $|\tau(\overline\chi)|=\sqrt q$. Theorem~\ref{thm:gauss}, applied to $\overline\chi$, gives
\[
\chi(n)=\frac1{\tau(\overline\chi)}\sum_{a=1}^{q-1}\overline\chi(a)e\Bigl(\frac{an}q\Bigr).
\]
The term $a=q$ is absent because $\overline\chi(q)=0$ for $q>1$. Insert the formula into $L(1,\chi)=\sum_n\chi(n)/n$. The sum over $a$ is finite and the series over $n$ converges, so they may be exchanged. Lemma~\ref{lem:logseries} at $x=2\pi a/q$ gives
\[
\begin{aligned}L(1,\chi)&=\frac1{\tau(\overline\chi)}\sum_{a=1}^{q-1}\overline\chi(a)\Bigl(-\log\Bigl(2\sin\frac{\pi a}q\Bigr)\eqbr+i\pi\Bigl(\frac12-\frac aq\Bigr)\Bigr).\end{aligned}
\]
The sum $\sum_a\overline\chi(a)$ is $0$, so the constant terms $\log 2$ and $i\pi/2$ drop out.

If $\chi$ is odd, then $\overline\chi$ is odd. The substitution $a\mapsto q-a$ sends $\log\sin(\pi a/q)$ to itself and sends $\overline\chi(a)$ to $-\overline\chi(a)$. The weighted sum of the logarithms equals its own negative, so it is $0$. The first formula remains.

If $\chi$ is even, the same substitution sends $\overline\chi(a)$ to itself and sends $a$ to $q-a$. Then $\sum_a\overline\chi(a)\,a=\sum_a\overline\chi(a)(q-a)$. The sum of $\overline\chi(a)$ vanishes, so $\sum_a\overline\chi(a)\,a$ equals its own negative and is $0$. The second formula remains.
\end{proof}

\begin{example}\label{ex:leibniz}
For $\chi_{-4}$: $\overline\chi=\chi$, $\tau(\chi_{-4})=e(\frac14)-e(\frac34)=2i$, and $\sum_a\chi(a)a=1-3=-2$, so $L(1,\chi_{-4})=-\frac{i\pi}{4\cdot2i}(-2)=\frac\pi4$, Leibniz's series.
\end{example}

For a real primitive character, $\tau(\chi)^2=\chi(-1)q$ (Exercise~\ref{ex:gausssquare}), so $\tau(\chi)$ is $\pm\sqrt q$ or $\pm i\sqrt q$. Gauss's theorem fixes the sign for the Legendre symbol: $\tau=\sqrt p$ if $p\equiv1\pmod4$, and $\tau=i\sqrt p$ if $p\equiv3\pmod4$ (Volume~5, Theorem~8.14). For $p\equiv3\pmod4$ the first formula becomes
\[
L(1,\chi)=-\frac\pi{p^{3/2}}\sum_{a<p}\Bigl(\frac ap\Bigr)a.
\]
Since $L(1,\chi)>0$ (Exercise~\ref{ex:lpositive}), the sum of $\bigl(\frac ap\bigr)a$ is negative. That is an arithmetic statement. No proof that avoids analysis is known (Exercise~\ref{ex:residuesum}).

\begin{remark}[for culture]
For $p\equiv3\pmod4$, $p>3$, the number $h(-p)=-\frac1p\sum_{a<p}\bigl(\frac ap\bigr)a$ is a positive integer, the \emph{class number} of $\Q(\sqrt{-p})$. It measures the failure of unique factorisation in the ring of integers (Volume~13, Chapter~2). Then $L(1,\chi)=\pi h(-p)/\sqrt p$, which is Dirichlet's class number formula, and
\[
\sum_{n<p/2}\Bigl(\frac np\Bigr)=\Bigl(2-\Bigl(\frac2p\Bigr)\Bigr)h(-p).
\]
That sum is the maximum of the walk in Figure~\ref{fig:charsums}. For even characters the second formula produces logarithms of units of real quadratic fields. For $q=5$ the unit contributes $\log\frac{1+\sqrt5}2$. These statements are quoted here. Proofs are in Davenport, Chapter~6, and the class number formula is proved in Volume~13, Chapter~11.
\end{remark}

\section{What we proved, and what comes next}

The indicator of a residue class is an average of characters (Proposition~\ref{prop:orth}). Every character is induced by exactly one primitive character, and the modulus of that primitive character is the conductor (Proposition~\ref{prop:conductor}). For a primitive character modulo $q$, the Gauss sum has modulus $\sqrt q$ (Theorem~\ref{thm:gauss}), and the same identity writes $\chi(n)$ as a short sum of exponentials. The series $L(s,\chi)$ has an Euler product for $\sigma>1$. If $\chi\ne\chi_0$, the series still converges for $\sigma>0$ (Theorem~\ref{thm:lbasic}). The fact the primes need is Theorem~\ref{thm:lnonzero}: $L(1,\chi)\ne0$ whenever $\chi\ne\chi_0$. Theorem~\ref{thm:l1formula} gives the value in finite terms. For $\chi_{-4}$ that value is $\pi/4$.

The class number formula in the remark was quoted, not proved. It needs the arithmetic of quadratic fields, and it waits for Volume~13.

The think box asked what $L(1,\chi_{-4})\ne0$ says about primes. Chapter~\ref{ch:dirichlet} answers the question for every modulus at once. The principal character is the only one whose weighted prime sum diverges. Every class prime to $q$ receives an equal share of that divergence.

\section{Exercises}

\begin{exercise}[$\star$]\label{ex:chartables}
Write the tables of the characters modulo $5$, $8$ and $12$. Which are real, which are odd, and which are primitive? Give the conductor of
each.
\end{exercise}

\begin{exercise}[$\star\star$]\label{ex:primcount}
Let $\varphi^*(q)$ be the number of primitive characters modulo $q$. Show that $\sum_{d\mid q}\varphi^*(d)=\varphi(q)$, deduce that $\varphi^*=\mu*\varphi$ is multiplicative
with $\varphi^*(p)=p-2$ and $\varphi^*(p^k)=p^k(1-\frac1p)^2$ for $k\ge2$, and that there are no primitive characters modulo $q\equiv2\pmod4$.
\end{exercise}

\begin{exercise}[$\star\star$]\label{ex:gausssquare}
Show that $\overline{\tau(\chi)}=\chi(-1)\tau(\overline\chi)$ for every character, and that $\tau(\chi)\tau(\overline\chi)=\chi(-1)q$ when $\chi$ is primitive. Deduce that
$\tau(\chi)^2=\chi(-1)q$ for a real primitive character.
\end{exercise}

\begin{exercise}[$\star$]\label{ex:lpositive}
Show that $L(1,\chi)>0$ for every real character $\chi\ne\chi_0$.
\end{exercise}

\begin{exercise}[$\star\star$]\label{ex:l1values}
\emph{Signature exercise.} Show that $L(1,\chi_{-3})=1-\frac12+\frac14-\frac15+\frac17-\cdots=\frac\pi{3\sqrt3}$ and
$L(1,\chi_5)=1-\frac12-\frac13+\frac14+\frac16-\cdots=\frac2{\sqrt5}\log\frac{1+\sqrt5}2$.
\end{exercise}

\begin{exercise}[$\star\star$]\label{ex:residuesum}
Let $p\equiv3\pmod4$ be a prime. Show that the sum of the quadratic non-residues in $\{1,\dots,p-1\}$ exceeds the sum of the quadratic
residues. Check it for $p=7$ and $p=11$.
\end{exercise}

\begin{exercise}[$\star\star$]\label{ex:lderivative}
Let $q\ge3$ and $\chi\ne\chi_0$ modulo $q$. Show that $L'(s,\chi)=-\sum_n\chi(n)(\log n)n^{-s}$ on $\sigma>0$, that $|L(1,\chi)|\le\log q+2$, and that
$|L'(1,\chi)|\le\frac12\log^2q+\log q+1$. (Split the series at $n=q$ and use Abel summation on the tail.)
\end{exercise}

\begin{exercise}[Journey]\label{pb:primroot}
\emph{Primitive roots as a sum of characters.} Let $p$ be an odd prime.
\begin{enumerate}[label=(\alph*)]
\item Show that for every $n$ prime to $p$,
\[
\mathbf 1_{\text{$n$ is a primitive root}}=\sum_{e\mid p-1}\frac{\mu(e)}e\sum_{\chi^e=\chi_0}\chi(n),
\]
the inner sum running over the $e$ characters modulo $p$ whose $e$-th power is the principal character $\chi_0$.
\item Sum (a) over $1\le n\le p-1$ and recover the number $\varphi(p-1)$ of primitive roots modulo $p$. Which characters survive the sum?
\item For $p=7$, write the indicator of (a) in terms of the six characters modulo $7$, and check it at $n=3$ and at $n=2$.
\end{enumerate}
The exercise stops at this character formula. It writes the indicator of the primitive roots as a sum of characters. Bounds for those sums belong to Chapter~\ref{ch:sums}.
\end{exercise}

\begin{exercise}[$\star$]\label{ex:chi4partial}
Compute $\sum_{n=1}^{12}\chi_{-4}(n)$ and the reduced fraction $\sum_{n=1}^{12}\chi_{-4}(n)/n$. Give the second sum as a decimal to four places, and its distance below $\pi/4=0.7854$ to four places.
\end{exercise}

\begin{exercise}[$\star\star$]\label{ex:mod5cols}
In Figure~\ref{fig:mod5grid}, add each column. Which columns sum to $4$, and which sum to $0$? For $n$ prime to $5$, deduce the same sums from Proposition~\ref{prop:orth}(b).
\end{exercise}

\begin{exercise}[$\star\star$]\label{ex:gauss12}
Let $\chi$ be the character modulo $12$ induced by $\chi_{-4}$, so $\chi(n)=\chi_{-4}(n)$ when $\gcd(n,12)=1$ and $\chi(n)=0$ otherwise. Compute $\tau(\chi)$ exactly. Compare $|\tau(\chi)|$ with $\sqrt{12}$, and name the hypothesis of Theorem~\ref{thm:gauss} that fails for $\chi$.
\end{exercise}

\begin{exercise}[$\star\star\star$]\label{ex:l2mod4}
Let $\chi_0$ be the principal character modulo $4$. Using Theorem~\ref{thm:lbasic}(c) and $\zeta(2)=\pi^2/6$ (Volume~5, Theorem~3.7), show that $L(2,\chi_0)=\pi^2/8$. Then compute $1+\frac19+\frac1{25}+\frac1{49}+\frac1{81}$ as a decimal to four places, and its distance below $\pi^2/8=1.2337$.
\end{exercise}

\section{Hints}

\begin{hint}{ex:chartables}
$(\Z/8\Z)^*\cong(\Z/2\Z)^2$ and $(\Z/12\Z)^*\cong(\Z/2\Z)^2$; a character is fixed by its values at generators. Use Proposition~\ref{prop:conductor} to test
the proper divisors $6$, $4$, $3$, $2$, $1$.
\end{hint}
\begin{hint}{ex:primcount}
Each character modulo $q$ is induced by exactly one primitive character modulo some $d\mid q$, and conversely each primitive character
modulo $d\mid q$ induces exactly one character modulo $q$. Then use Möbius inversion.
\end{hint}
\begin{hint}{ex:gausssquare}
Conjugate the definition and substitute $a\mapsto-a$. Then combine with $|\tau(\chi)|^2=q$.
\end{hint}
\begin{hint}{ex:lpositive}
$L(\sigma,\chi)>0$ for real $\sigma>1$, and $L(\sigma,\chi)$ is continuous on $\left]0,\infty\right[$.
\end{hint}
\begin{hint}{ex:l1values}
Theorem~\ref{thm:l1formula}: compute $\tau(\chi_{-3})$ and $\tau(\chi_5)$ directly, and $\sin\frac{2\pi}5=2\sin\frac\pi5\cos\frac\pi5$ with $2\cos\frac\pi5=\frac{1+\sqrt5}2$.
\end{hint}
\begin{hint}{ex:residuesum}
Combine the first formula of Theorem~\ref{thm:l1formula} with Gauss's sign and Exercise~\ref{ex:lpositive}.
\end{hint}
\begin{hint}{ex:lderivative}
For $n\le q$ bound each term trivially. For $n>q$, write the tail with $A(t)=\sum_{n\le t}\chi(n)$, $|A(t)|\le\varphi(q)\le q$.
\end{hint}

\begin{hint}{pb:primroot}
(a) $\sum_{\chi^e=\chi_0}\chi(n)$ is $e$ or $0$ according as $n$ is an $e$-th power modulo $p$ or not. (b) A non-principal character sums to $0$ over a full period.
(c) The characters with $\chi^e=\chi_0$ are those of order dividing $e$, for $e\in\{1,2,3,6\}$.
\end{hint}

\begin{hint}{ex:chi4partial}
The non-zero terms are $1-\frac13+\frac15-\frac17+\frac19-\frac1{11}$. The common denominator is $3465$.
\end{hint}
\begin{hint}{ex:mod5cols}
Add $1+i-1-i$ in the column $n=2$. For the proposition, take $a=1$, and use $\overline{\chi(1)}=1$.
\end{hint}
\begin{hint}{ex:gauss12}
The surviving residues are $1,5,7,11$, with values $1,1,-1,-1$. Use $e(1/12)=\frac{\sqrt3}2+\frac i2$ and the angles $150^\circ$, $210^\circ$, $330^\circ$.
\end{hint}
\begin{hint}{ex:l2mod4}
The factor at the prime $2$ is $1-2^{-2}=\frac34$. The five-term sum has denominator $3^4\cdot5^2\cdot7^2=99225$.
\end{hint}

\section{Solutions}

\begin{solution}{ex:chartables}
\idea A character of a product of cyclic groups of exponent $2$ is a choice of $\pm1$ at each generator. The conductor is the least $d$ such that the character is $1$ on every unit congruent to $1$ modulo $d$.
\writeup Modulo $5$, with generator $2$, the value $\chi(2)$ runs through $\{1,i,-1,-i\}$. These give $\chi_0$, two complex odd characters (because $\chi(4)=\chi(2)^2=-1$ and $4\equiv-1$), and the Legendre character $\chi_5$, which is real and even. The modulus is prime, so the three non-principal characters are primitive. Modulo $8$, the unit group is $\{1,3,5,7\}$. Every element has order $2$, and $3$ and $5$ generate the group. The four characters are real. They are $\chi_0$; the character $\chi_{-4}$ with values $1,-1,1,-1$ on $1,3,5,7$, of conductor $4$, odd; the character $\chi_8$ with values $1,-1,-1,1$, so that $\chi_8(n)=1$ exactly for $n\equiv\pm1\pmod8$, of conductor $8$, even; and $\chi_{-8}=\chi_{-4}\chi_8$ with values $1,1,-1,-1$, of conductor $8$, odd. Modulo $12$, the unit group is $\{1,5,7,11\}$. The characters are $\chi_0$; the character induced by $\chi_{-4}$, with values $1,1,-1,-1$, of conductor $4$, odd; the character induced by $\chi_{-3}$, with values $1,-1,1,-1$, of conductor $3$, odd; and their product $\chi_{12}$, with values $1,-1,-1,1$. The product is not trivial on $U_6=U_3=\{1,7\}$, nor on $U_4=\{1,5\}$, nor on $U_2$, nor on $U_1$. It is primitive and even, and $\chi_{12}(n)=1$ exactly when $n\equiv\pm1\pmod{12}$.
\end{solution}

\begin{solution}{ex:primcount}
\idea Group the characters modulo $q$ by conductor. Each conductor $d$ contributes its primitive characters, and no others.
\writeup By Proposition~\ref{prop:conductor}(c), the map $\chi\mapsto(d^*,\chi^*)$ sends the characters modulo $q$ to pairs consisting of a divisor $d\mid q$ and a primitive character modulo $d$. The map is injective: $\chi^*$ determines $\chi$, because the two agree on integers prime to $q$ and $\chi$ vanishes elsewhere. The map is surjective: a primitive character $\psi$ modulo $d\mid q$ induces $\chi=\chi_0\psi$ modulo $q$, and uniqueness in (c) forces $\chi^*=\psi$. Counting gives $\varphi(q)=\sum_{d\mid q}\varphi^*(d)$. Möbius inversion (Volume~2, Theorem~4.4) gives $\varphi^*=\mu*\varphi$. A convolution of multiplicative functions is multiplicative. At a prime power, $\varphi^*(p)=\varphi(p)-\varphi(1)=p-2$, and for $k\ge2$
\[
\varphi^*(p^k)=\varphi(p^k)-\varphi(p^{k-1})=p^{k-2}(p-1)^2=p^k\Bigl(1-\frac1p\Bigr)^2.
\]
In particular $\varphi^*(2)=0$. If $q=2m$ with $m$ odd, then $\varphi^*(q)=\varphi^*(2)\varphi^*(m)=0$. Modulo $12$ the same formula gives $\varphi^*(4)\varphi^*(3)=1\cdot1=1$, which matches the one primitive character found in Exercise~\ref{ex:chartables}.
\end{solution}

\begin{solution}{ex:gausssquare}
\idea Conjugation sends $e(a/q)$ to $e(-a/q)$ and $\chi(a)$ to $\overline\chi(a)$. The reflection $a\mapsto-a$ turns the conjugated sum back into a Gauss sum, and the price is the factor $\chi(-1)$.
\writeup Substitute $b=-a$:
\[
\begin{aligned}\overline{\tau(\chi)}&=\sum_a\overline\chi(a)e\Bigl(-\frac aq\Bigr)=\sum_b\overline\chi(-b)e\Bigl(\frac bq\Bigr)\eqbr=\overline\chi(-1)\tau(\overline\chi).\end{aligned}
\]
Since $\chi(-1)=\pm1$, one has $\overline\chi(-1)=\chi(-1)$. Theorem~\ref{thm:gauss} then gives
\[
\begin{aligned}\tau(\chi)\tau(\overline\chi)&=\chi(-1)\tau(\chi)\overline{\tau(\chi)}\eqbr=\chi(-1)|\tau(\chi)|^2=\chi(-1)q.\end{aligned}
\]
If $\chi$ is real, then $\overline\chi=\chi$, so $\tau(\chi)^2=\chi(-1)q$. Check the sign on the character from the opening: $\tau(\chi_{-4})^2=(2i)^2=-4$, and $\chi_{-4}(-1)\cdot4=-4$.
\end{solution}

\begin{solution}{ex:lpositive}
\idea On the real axis past $1$, the series $L(\sigma,\chi)$ is a positive real exponential. A continuous function that stays positive on $(1,\infty)$ can reach a negative value at $\sigma=1$ only by passing through $0$, and Theorem~\ref{thm:lnonzero} forbids that zero.
\writeup For $\sigma>1$, Theorem~\ref{thm:lbasic}(a) gives $L(\sigma,\chi)=\exp(S)$ with $S$ real, because $\chi$ is real, so $L(\sigma,\chi)>0$. By Theorem~\ref{thm:lbasic}(b) the map $\sigma\mapsto L(\sigma,\chi)$ is continuous on $\left]0,\infty\right[$. Therefore $L(1,\chi)=\lim_{\sigma\to1^+}L(\sigma,\chi)\ge0$. Theorem~\ref{thm:lnonzero} says the value is not $0$.
\end{solution}

\begin{solution}{ex:l1values}
\idea Both characters are real and primitive. Compute the Gauss sum by pairing terms, then evaluate the short finite sum in Theorem~\ref{thm:l1formula}.
\writeup The character $\chi_{-3}$ is odd and primitive, and
\[
\tau(\chi_{-3})=e\Bigl(\frac13\Bigr)-e\Bigl(\frac23\Bigr)=2i\sin\frac{2\pi}3=i\sqrt3.
\]
The first formula of Theorem~\ref{thm:l1formula} gives
\[
L(1,\chi_{-3})=-\frac{i\pi}{3\cdot i\sqrt3}(1-2)=\frac\pi{3\sqrt3}=0.6046\ldots.
\]
The character $\chi_5$ is even and primitive. Pairing $a$ with $5-a$,
\[
\tau(\chi_5)=2\cos\frac{2\pi}5-2\cos\frac{4\pi}5=\frac{\sqrt5-1}2+\frac{\sqrt5+1}2=\sqrt5.
\]
Since $\sin(4\pi/5)=\sin(\pi/5)$ and $\sin(3\pi/5)=\sin(2\pi/5)$, the second formula gives
\[
\begin{aligned}L(1,\chi_5)&=-\frac1{\sqrt5}\Bigl(2\log\sin\frac\pi5-2\log\sin\frac{2\pi}5\Bigr)\eqbr=\frac2{\sqrt5}\log\frac{\sin(2\pi/5)}{\sin(\pi/5)}.\end{aligned}
\]
The double-angle formula gives $\sin(2\pi/5)/\sin(\pi/5)=2\cos(\pi/5)=(1+\sqrt5)/2$, the golden ratio. Therefore
\[
L(1,\chi_5)=\frac2{\sqrt5}\log\frac{1+\sqrt5}2=0.4304\ldots.
\]
Both values agree with the partial sums of the two series, which checks the signs. The golden ratio is the fundamental unit of $\Z\bigl[\frac{1+\sqrt5}2\bigr]$. That is the even case of Dirichlet's class number formula, quoted in the remark above.
\end{solution}

\begin{solution}{ex:residuesum}
\idea For an odd real character, Dirichlet's formula writes the positive number $L(1,\chi)$ as a negative multiple of $\sum\chi(a)a$, once Gauss's sign is known. The sum $\sum\chi(a)a$ is therefore negative, and that is the comparison of residues with non-residues.
\writeup Let $\chi=\bigl(\frac\cdot p\bigr)$. The character is odd because $p\equiv3\pmod4$, and it is primitive. By Volume~5, Theorem~8.14, $\tau(\chi)=i\sqrt p$. The first formula of Theorem~\ref{thm:l1formula} gives
\[
L(1,\chi)=-\frac{i\pi}{p\cdot i\sqrt p}\sum_a\chi(a)a=-\frac\pi{p^{3/2}}\sum_a\chi(a)a.
\]
Exercise~\ref{ex:lpositive} says $L(1,\chi)>0$, so $\sum_a\chi(a)a<0$. The non-residues add to more than the residues. For $p=7$ the residues are $1,2,4$, with sum $7$, and the non-residues are $3,5,6$, with sum $14$. For $p=11$ the residues are $1,3,4,5,9$, with sum $22$, and the non-residues are $2,6,7,8,10$, with sum $33$. In both cases the difference of the two sums, divided by $p$, equals $1$. That integer is $h(-7)=h(-11)=1$. No proof that avoids analysis is known.
\end{solution}

\begin{solution}{ex:lderivative}
\idea The first $q$ terms are at most a harmonic sum. The tail is an Abel sum against partial sums bounded by $\varphi(q)$.
\writeup The proof of Volume~5, Proposition~10.1(b) gives uniform convergence of the partial sums of $L(s,\chi)$ on closed discs in $\sigma>0$. Volume~5, Theorem~9.15(a) allows term-by-term differentiation, so $L'(s,\chi)=-\sum_n\chi(n)(\log n)n^{-s}$ on $\sigma>0$. Let $A(t)=\sum_{n\le t}\chi(n)$. Then $|A(t)|\le\varphi(q)\le q$, and $A(q)=0$ by Proposition~\ref{prop:orth}(a). For a $C^1$ function $f$ that decreases to $0$ on $[q,\infty[$, Lemma~\ref{lem:abelint} gives
\[
\sum_{q<n\le x}\chi(n)f(n)=A(x)f(x)-\int_q^x A(t)f'(t)\,\dd t.
\]
The tail $\sum_{n>q}\chi(n)f(n)$ therefore has absolute value at most $\varphi(q)\int_q^\infty(-f')=\varphi(q)f(q)$.

With $f(t)=1/t$, the head satisfies $\bigl|\sum_{n<q}\chi(n)/n\bigr|\le1+\log q$, by comparison with $\int_1^q dt/t$ (Volume~5, Corollary~11.6). The tail is at most $\varphi(q)/q\le1$. Hence $|L(1,\chi)|\le\log q+2$.

With $f(t)=(\log t)/t$, which decreases on $[e,\infty[$ and so on $[q,\infty[$, the tail is at most $\varphi(q)(\log q)/q\le\log q$. For the head, $(\log n)/n\le\int_{n-1}^n(\log t)/t\,\dd t$ when $n\ge4$, so
\[
\sum_{n<q}\frac{\log n}n\le\frac{\log2}2+\frac{\log3}3+\int_3^q\frac{\log t}t\,\dd t\le1+\frac12\log^2 q.
\]
Adding the head and the tail gives the bound on $|L'(1,\chi)|$. The bounds are crude, and they are uniform in $\chi$. Chapter~\ref{ch:dirichlet} needs only that $L(1,\chi)$ and $L'(1,\chi)$ are finite. Chapter~\ref{ch:sums} sharpens estimates of this kind.
\end{solution}

\begin{solution}{pb:primroot}
\idea A primitive root is a unit whose index is coprime to $p-1$. Möbius inversion over the divisors of $p-1$ writes that condition as a sum of characters.

\writeup (a) The characters with $\chi^e=\chi_0$ form a group of order $e$. They are the characters of $(\Z/p\Z)^*$ trivial on the subgroup of $e$-th powers. By orthogonality, their sum at $n$ equals $e$ if $n$ is an $e$-th power modulo $p$, and equals $0$ otherwise. Write $n=g^k$ for a primitive root $g$. Then $n$ is an $e$-th power if and only if $e\mid k$, using $e\mid p-1$. The right side of the formula is therefore $\sum_{e\mid\gcd(k,p-1)}\mu(e)$. That sum equals $1$ when $\gcd(k,p-1)=1$, which is the condition that $n$ be a primitive root, and equals $0$ otherwise.

(b) For $\chi\ne\chi_0$, the sum $\sum_{n=1}^{p-1}\chi(n)$ is $0$, by Proposition~\ref{prop:orth}(a). For each $e$, exactly one of the $e$ characters with $\chi^e=\chi_0$ is principal. Only those principal characters survive. Each contributes $p-1$, and the number of primitive roots is
\[
(p-1)\sum_{e\mid p-1}\frac{\mu(e)}e=\varphi(p-1).
\]
The last step is $\varphi(m)/m=\sum_{e\mid m}\mu(e)/e$ (Volume~2, Chapter~4).

(c) Here $p-1=6$. The squarefree divisors are $e=1,2,3,6$, with $\mu(e)=1,-1,-1,1$. Let $\omega=e(1/6)$, and let $\chi$ be the character with $\chi(3)=\omega$. The residue $3$ is a primitive root modulo $7$. The six characters are $\chi^j$ for $0\le j\le5$, and $\chi^j$ has order dividing $e$ exactly when $6/e$ divides $j$. At $n=3$ the inner sums are $1$ for $e=1$, then $1+\omega^3=0$ for $e=2$, then $1+\omega^2+\omega^4=0$ for $e=3$, and $\sum_{j=0}^5\omega^j=0$ for $e=6$. The indicator equals $1$, and $3$ is a primitive root. At $n=2=3^2$ one has $\chi^j(2)=\omega^{2j}$. The inner sums are $1$, then $1+\omega^6=2$, then $1+\omega^4+\omega^8=0$, then $\sum_{j=0}^5\omega^{2j}=0$. The indicator equals $1-\frac12\cdot2=0$. Indeed $2^3=8\equiv1\pmod7$, so $2$ has order $3$ and does not generate $(\Z/7\Z)^*$.

\tools{Möbius inversion (Volume~2, Chapter~4); orthogonality of characters (Volume~2, Theorem~11.5, and Proposition~\ref{prop:orth}).}

\leads The same expansion writes the indicator of any subgroup condition on $(\Z/p\Z)^*$. This exercise stops once the indicator is a sum of characters.
\end{solution}

\begin{solution}{ex:chi4partial}
\idea Three full periods of a non-principal character sum to $0$. The weighted sum keeps six non-zero terms, one for each odd integer up to $11$.
\writeup The values of $\chi_{-4}$ on $1,\dots,12$ are $1,0,-1,0,1,0,-1,0,1,0,-1,0$. Their sum is $0$. The weighted sum is
\[
1-\frac13+\frac15-\frac17+\frac19-\frac1{11}.
\]
The least common denominator of $1,3,5,7,9,11$ is $3^2\cdot5\cdot7\cdot11=3465$. The sum equals
\[
\frac{3465-1155+693-495+385-315}{3465}=\frac{2578}{3465}.
\]
The fraction is in lowest terms. Dividing gives $0.7440115\ldots$, which is $0.7440$ to four decimal places. Since $\pi/4=0.785398\ldots$ is $0.7854$ to four places, the partial sum lies $0.0414$ below $\pi/4$.
\end{solution}

\begin{solution}{ex:mod5cols}
\idea The column $n=1$ is the only column of ones. In every other column the real values cancel and the imaginary values cancel.
\writeup Reading Figure~\ref{fig:mod5grid} from $\chi_0$ down to $\overline\psi$: the column $n=0$ is $0+0+0+0=0$; the column $n=1$ is $1+1+1+1=4$; the column $n=2$ is $1+i-1-i=0$; the column $n=3$ is $1-i-1+i=0$; the column $n=4$ is $1-1+1-1=0$. For $n$ prime to $5$, Proposition~\ref{prop:orth}(b) with $a=1$ says
\[
\frac1{\varphi(5)}\sum_\chi\chi(n)=1
\]
if $n\equiv1\pmod5$, and equals $0$ otherwise, because $\overline{\chi(1)}=1$ for every character. Here $\varphi(5)=4$, so the column sums are $4$ at $n\equiv1$ and $0$ at the other units. At multiples of $5$ every character is $0$, so that column sums to $0$ as well.
\end{solution}

\begin{solution}{ex:gauss12}
\idea Theorem~\ref{thm:gauss} needs a primitive character. This $\chi$ is $\chi_{-4}$ with the multiples of $3$ removed, so the Gauss sum is a sum of four exponentials, and its modulus need not be $\sqrt{12}$.
\writeup The units modulo $12$ are $1,5,7,11$, and $\chi$ takes the values $1,1,-1,-1$ there, since those residues are $1,1,3,3$ modulo $4$. Thus
\[
\tau(\chi)=e\Bigl(\frac1{12}\Bigr)+e\Bigl(\frac5{12}\Bigr)-e\Bigl(\frac7{12}\Bigr)-e\Bigl(\frac{11}{12}\Bigr).
\]
The four exponentials are $\frac{\sqrt3}2+\frac i2$, $-\frac{\sqrt3}2+\frac i2$, $-\frac{\sqrt3}2-\frac i2$ and $\frac{\sqrt3}2-\frac i2$. Adding with the signs $+,+,-,-$ cancels every real part and leaves $i/2+i/2+i/2+i/2=2i$. So $|\tau(\chi)|=2$, while $\sqrt{12}=2\sqrt3\approx3.464$. The character is trivial on $U_4=\{1,5\}$, and $4$ is a proper divisor of $12$, so the conductor is $4$, not $12$. Theorem~\ref{thm:gauss} assumes primitivity, and $\chi$ is not primitive.
\end{solution}

\begin{solution}{ex:l2mod4}
\idea The principal character modulo $4$ deletes the Euler factor at $2$. At $s=2$ that factor is $3/4$, and $\zeta(2)$ is already known.
\writeup Theorem~\ref{thm:lbasic}(c) gives $L(s,\chi_0)=\zeta(s)(1-2^{-s})$ for $\sigma>1$. At $s=2$, Volume~5, Theorem~3.7 supplies $\zeta(2)=\pi^2/6$, and $1-1/4=3/4$, so
\[
L(2,\chi_0)=\frac{\pi^2}6\cdot\frac34=\frac{\pi^2}8.
\]
The same value is the sum of the reciprocals of the odd squares. The first five terms are
\[
1+\frac19+\frac1{25}+\frac1{49}+\frac1{81}.
\]
The common denominator is $3^4\cdot5^2\cdot7^2=99225$, and the sum is $117469/99225=1.1838649\ldots$, which is $1.1839$ to four decimal places. Since $\pi^2/8=1.233700\ldots$ is $1.2337$ to four places, the partial sum lies $0.0498$ below $\pi^2/8$.
\end{solution}

\section*{Sources and further reading}

\begin{itemize}
\item P.~G.~L.~Dirichlet, \emph{Beweis des Satzes, dass jede unbegrenzte arithmetische Progression, deren erstes Glied und Differenz ganze
Zahlen ohne gemeinschaftlichen Factor sind, unendlich viele Primzahlen enthält}, Abhandlungen der Königlichen Preussischen Akademie
der Wissenschaften (1837), 45--81. Read to the Academy on 27 July 1837.
\item H.~Davenport, \emph{Multiplicative Number Theory}, 3rd ed., revised by H.~L.~Montgomery, Springer GTM~74, 2000, Chapters~1, 2 and~4--6
(characters, Gauss sums, $L(1,\chi)\ne0$, primitive characters, the class number formula).
\item T.~M.~Apostol, \emph{Introduction to Analytic Number Theory}, Springer, 1976, Chapters~6--8 (characters, Dirichlet's theorem, Gauss sums, induced
characters).
\item J.-P.~Serre, \emph{A Course in Arithmetic}, Springer GTM~7, 1973, Chapter~VI (Dirichlet's theorem with Landau's lemma on series with
positive coefficients, the route followed here).
\item J.~Elstrodt, \emph{The life and work of Gustav Lejeune Dirichlet (1805--1859)}, Clay Mathematics Proceedings~7 (2007), 1--37. Elstrodt
records Jacobi's letter of 10 September 1838 to Moritz Hermann Jacobi. The English sentence in the epigraph is Elstrodt's translation.
\end{itemize}
